\documentclass[11pt]{article}
\usepackage[margin=1in]{geometry}
\usepackage{fontspec}
\usepackage{setspace}
\usepackage{hyperref}
\hypersetup{colorlinks=true, allcolors=black}
\setlength{\parindent}{0pt}
\setlength{\parskip}{0.5em}
\setstretch{1.08}
\pagestyle{empty}
\setmainfont{Times New Roman}

\begin{document}

\begin{flushright}
Ke Cheng HSU (Clio)\\
Department of Economics, HKUST\\
Clear Water Bay, Kowloon, Hong Kong\\
+852 6840 9753\\
\href{mailto:khsu@connect.ust.hk}{khsu@connect.ust.hk}\\
\href{https://cliohsu1997.github.io}{cliohsu1997.github.io}
\end{flushright}

2 November 2026

Faculty Search Committee\\
International College of Liberal Arts\\
Yamanashi Gakuin University\\
\href{mailto:faculty.search@c2c.ac.jp}{faculty.search@c2c.ac.jp}

Dear Members of the Faculty Search Committee,

I am writing to apply for the Lecturer / Associate Professor / Professor position in Economics and Finance at the International College of Liberal Arts (JREC-IN D126101270). I am a Ph.D.\ candidate in Economics at the Hong Kong University of Science and Technology. My expected graduation date is June 2027, and I can move my defense forward and complete the degree earlier upon receiving an offer so that I can begin on 1 April 2027.

iCLA's English-medium liberal arts program is why I am applying. The post asks for undergraduate courses in macroeconomics, intermediate macroeconomics, and econometrics, and welcomes international finance, risk management, and portfolio management. From 2020 to 2025 I was a teaching assistant for undergraduate macroeconomics, the economics of uncertainty and information, and Ph.D.\ microeconomic theory, with an average evaluation of 4.72/5. In 2026 I designed and taught the mathematics camp for the MSc in Economics (4.53/5). For portfolio management and risk management, my undergraduate degree is in Mathematics and Economics with a minor in Actuarial Mathematics, so I am already familiar with mean--variance portfolio choice, the CAPM, diversification, hedging, loss distributions, and value at risk. I would be glad to mentor students, develop courses for the Global Business and Economics program, and contribute to service at the college.

My research is empirical. The job-market paper estimates how an online return policy changes purchases and returns. A second paper, with Pak Hung Au, studies how firms use product design to direct the search behavior of consumers with different search costs. Both are attached. With Kohei Kawaguchi and Kei Kawai, I study how a price commitment at Japanese bookstores leaves unsold copies when demand is low. With Ce Bian and Liang Zhao, I study how product returns change the pass-through of a tax and who bears it.

I attach the iCLA CV, teaching philosophy, two papers, course evaluations, and signed confirmation of expected graduation. Pak Hung Au, Kohei Kawaguchi, and Kei Kawai will email reference letters directly. You can reach me at khsu@connect.ust.hk or +852 6840 9753.

Sincerely,

Ke Cheng HSU

\end{document}
